\section{Materials and Methods}
\label{sec:methods}

This section defines the graph-classification task, the Matrix-Residual architecture family, the benchmark protocol, and the mechanism metrics. The guiding principle is to isolate residual topology while holding the remaining learning pipeline fixed.

\subsection{Problem definition and notation}

Let \(\mathcal{G}=(\mathcal{V},\mathcal{E})\) be a graph with node feature matrix \(\mathbf{X}\in\mathbb{R}^{n\times d_{\mathrm{in}}}\) and adjacency structure \(\mathbf{A}\). For a dataset \(\mathcal{D}=\{(\mathcal{G}_i,y_i)\}_{i=1}^{N}\), the objective is to learn \(f_{\theta}:\mathcal{G}\rightarrow \hat{y}\), where \(\hat{y}\) is the predicted graph label.

The compared model families share the same graph-classification wrapper:
\begin{itemize}
    \item \textbf{Plain}: stacked message passing without explicit residual reuse.
    \item \textbf{VerticalRes}: same-branch residual reuse from the previous layer.
    \item \textbf{HorizontalRes}: cross-branch exchange between adjacent branches at the same depth.
    \item \textbf{MatrixRes}: local two-dimensional reuse over vertical, horizontal, and diagonal neighbors.
    \item \textbf{MatrixResGated}: matrix reuse with gated or sparse residual filtering.
\end{itemize}

\subsection{Shared graph-classification pipeline}

Let \(B\) be the number of branches and \(L\) be the number of message-passing layers. The hidden state at branch \(b\) and layer \(\ell\) is written as \(\mathbf{H}^{(\ell)}_{b}\), where \(b\in\{1,\ldots,B\}\). A branch-local propagation step is
\begin{equation}
\mathbf{Z}^{(\ell)}_{b} = \phi^{(\ell)}_{b}\left(\mathbf{H}^{(\ell-1)}_{b}, \mathbf{A}\right),
\end{equation}
where \(\phi\) is instantiated as one of four standard operators: GCNConv \cite{kipf2017gcn}, GATConv \cite{velivckovic2017graph}, SAGEConv \cite{hamilton2017inductive}, or GINConv \cite{xu2019gin}. The residual-enhanced update is
\begin{equation}
\mathbf{H}^{(\ell)}_{b} =
\mathbf{Z}^{(\ell)}_{b} +
\sum_{(b',\ell')\in\mathcal{N}(b,\ell)}
\mathcal{R}_{(b',\ell')\rightarrow(b,\ell)}
\left(\mathbf{H}^{(\ell')}_{b'}\right).
\end{equation}
The only term that changes across architectures is the residual neighborhood \(\mathcal{N}(b,\ell)\) and the residual transform \(\mathcal{R}\).

After the final message-passing layer, each branch is pooled to a graph-level vector by global mean pooling and the branch vectors are combined before classification. We intentionally avoid architecture-specific pooling heads such as differentiable pooling or self-attention pooling \cite{ying2018hierarchical,lee2019self}, because the purpose is to isolate residual topology rather than to compare readout mechanisms. The classifier is trained with cross-entropy loss.

\subsection{Residual neighborhoods}
\label{sec:model}

\textbf{Plain.} The plain baseline has no residual neighborhood:
\begin{equation}
\mathcal{N}(b,\ell)=\emptyset.
\end{equation}

\textbf{VerticalRes.} Vertical residual reuse injects the previous layer on the same branch:
\begin{equation}
\mathcal{N}(b,\ell)=\{(b,\ell-1)\}.
\end{equation}

\textbf{HorizontalRes.} Horizontal residual reuse exchanges information across adjacent branches at the same depth:
\begin{equation}
\mathcal{N}(b,\ell)=\{(b-1,\ell),(b+1,\ell)\},
\end{equation}
with boundary branches using only valid neighbors.

\textbf{MatrixRes.} Matrix residual reuse combines vertical, horizontal, and diagonal neighborhoods:
\begin{equation}
\mathcal{N}(b,\ell)=
\{(b,\ell-1),(b-1,\ell),(b+1,\ell),(b-1,\ell-1),(b+1,\ell-1)\}.
\end{equation}
This local stencil is intentionally not dense all-to-all mixing. It preserves an interpretable residual structure over the branch-by-layer grid.

\textbf{MatrixResGated.} MatrixResGated uses the same neighborhood as MatrixRes, but the residual transform controls the injected signal:
\begin{equation}
\mathcal{R}(\mathbf{U}) =
\begin{cases}
\alpha \mathbf{U}, & \text{identity residual filtering with a gate},\\
\alpha\,\mathrm{softshrink}_{\lambda}(\mathbf{U}), & \text{sparse residual filtering with a gate}.
\end{cases}
\end{equation}
In the main benchmark, Plain, VerticalRes, HorizontalRes, and MatrixRes use identity residual filtering. MatrixResGated uses sparse residual filtering with \(\lambda=0.05\) and a learnable scalar gate initialized at 0.8. The implementation also supports top-\(k\) residual filtering and fixed gates, but these are not the default settings used for the main reported results.

\subsection{Datasets and evaluation protocol}
\label{sec:datasets}

The main benchmark uses six datasets from TUDataset \cite{morris2020tudataset}: PROTEINS, DD, ENZYMES, MUTAG, AIDS, and Mutagenicity. PROTEINS and DD are binary protein-structure graph classification benchmarks, ENZYMES is a six-class enzyme classification benchmark, and MUTAG, AIDS, and Mutagenicity provide supplementary molecular-graph robustness checks. All datasets are loaded through PyTorch Geometric \cite{fey2019fast}. All experiments use the raw benchmark features and graph structures without task-specific feature engineering, graph augmentation, or edge redesign.

All main results use graph-level five-fold evaluation. For fold \(k\), the test accuracy is
\begin{equation}
\mathrm{Acc}_{k}=
\frac{1}{|\mathcal{T}_{k}|}
\sum_{(\mathcal{G}_{i},y_{i})\in\mathcal{T}_{k}}
\mathbf{1}[\hat{y}_{i}=y_{i}].
\end{equation}
Across folds, we report
\begin{equation}
\overline{\mathrm{Acc}}=\frac{1}{5}\sum_{k=1}^{5}\mathrm{Acc}_{k},
\end{equation}
and the fold-level standard deviation.

\subsection{Training configuration}

All main benchmark runs use the default branch count \(B=3\). The benchmark is repeated under four message-passing operators: GCNConv, GATConv, SAGEConv, and GINConv. The hidden dimension is 64 unless a sensitivity setting changes it, and all model families use the same optimizer, early-stopping logic, graph pooling, and fold protocol. Models are optimized with Adam \cite{kingma2014adam}; the checkpoint with the lowest validation loss is retained for test evaluation. The branch-count ablation evaluates \(B=1,\ldots,8\) for HorizontalRes, MatrixRes, and MatrixResGated on PROTEINS and DD under GCNConv. The sensitivity scan evaluates MatrixRes and MatrixResGated on fold 0 at \(B=3\), varying learning rate, dropout, hidden dimension, sparsity strength, and gate initialization. Selected MatrixResGated candidates are then rerun across all five folds.

\begin{table}[h]
\centering
\small
\caption{Key experimental settings shared by the main benchmark.}
\label{tab:hyperparams}
\begin{tabular}{l c}
\toprule
Parameter & Value \\
\midrule
Message-passing operators & GCNConv, GATConv, SAGEConv, GINConv \\
Default branch count & 3 \\
Default hidden dimension & 64 \\
Evaluation & Five-fold graph classification \\
Branch-count ablation & \(B=1,\ldots,8\) \\
Primary metric & Mean best test accuracy \\
Sensitivity scope & Fold 0, \(B=3\) \\
Early stopping & Validation loss, patience 60 or 80 \\
\bottomrule
\end{tabular}
\end{table}

\subsection{Mechanism metrics}

To interpret branch-count behavior, we track three mechanism indicators. Branch diversity measures the mean pairwise distance between branch representations. Branch cosine similarity measures whether branches remain redundant. Mean gradient norm tracks optimization weakening as branch count and residual traffic grow. CKA and residual-traffic summaries are retained as supplementary diagnostics.
